\subsection{مبانی نظری و فرمول‌بندی ریاضی}
فرض کنید تعداد سرهای پرس‌وجو $H$ و تعداد سرهای کلید/مقدار $H_{\text{kv}}$ باشد به‌طوری‌که $H_{\text{kv}} \le H$ و نسبت تکرار $R = H / H_{\text{kv}}$ یک عدد صحیح مثبت باشد. تانسورهای پرس‌وجو، کلید و مقدار با ابعاد زیر استخراج می‌شوند:
\begin{equation}
\mathbf{Q} \in \mathbb{R}^{B \times H \times N_q \times d}, \quad \mathbf{K} \in \mathbb{R}^{B \times H_{\text{kv}} \times N_k \times d}, \quad \mathbf{V} \in \mathbb{R}^{B \times H_{\text{kv}} \times N_k \times d}
\end{equation}
نوآوری اصلی \lr{WGQA} در معرفی بردارهای وزنی یادگیرنده $\mathbf{w}_k \in \mathbb{R}^{H_{\text{kv}}}$ و $\mathbf{w}_v \in \mathbb{R}^{H_{\text{kv}}}$ نهفته است. پیش از عملیات تکرار سرها، مقیاس‌دهی با این ضرایب انجام می‌پذیرد:
\begin{equation}
\tilde{\mathbf{K}}_g = \mathbf{w}_{k, g} \cdot \mathbf{K}_g, \quad \tilde{\mathbf{V}}_g = \mathbf{w}_{v, g} \cdot \mathbf{V}_g \quad \forall g \in \{1, \dots, H_{\text{kv}}\}
\end{equation}
سپس تانسورهای وزن‌دار شده با استفاده از تابع تکرار کپی‌شده (\lr{repeat\_kv}) به تعداد سرهای پرس‌وجو بسط داده می‌شوند:
\begin{equation}
\hat{\mathbf{K}} = \operatorname{repeat\_kv}(\tilde{\mathbf{K}}, R), \quad \hat{\mathbf{V}} = \operatorname{repeat\_kv}(\tilde{\mathbf{V}}, R)
\end{equation}
در صورت تجهیز مدل به تعبیه‌های مکانی چرخشی (\lr{RoPE})، نگاشت زاویه‌ای بر روی جفت‌های متعامد پرس‌وجو و کلید اعمال می‌گردد:
\begin{equation}
\mathbf{Q}_{\text{rot}}, \hat{\mathbf{K}}_{\text{rot}} = \operatorname{apply\_rotary\_pos\_emb}(\mathbf{Q}, \hat{\mathbf{K}}, \cos, \sin)
\end{equation}
در نهایت، ماتریس امتیازات توجه مقیاس‌پذیر و توزیع وزن با تابع بیشینه هموار محاسبه شده و خروجی استخراج می‌گردد:
\begin{equation}
\mathbf{A} = \operatorname{softmax}\left(\frac{\mathbf{Q} \hat{\mathbf{K}}^\top}{\sqrt{d}} + \mathbf{M}\right), \quad \mathbf{O} = \mathbf{A} \hat{\mathbf{V}}
\end{equation}
که در آن $\mathbf{M}$ ماتریس ماسک توجه (شامل ماسک علّی در دیکودر یا ماسک طول معتبر در انکودر) است.


\paragraph{تحلیل طیفی و نامساوی‌های پایداری خطا}
برای بررسی میزان دقت تقریب ماتریس وزن یا نگاشت توجه در ماژول \lr{wgqa\_attention}، تفاضل ماتریس فاکتورگیری‌شده از نگاشت متراکم اصلی بر اساس نامساوی هافدینگ-فون‌نیومان ارزیابی می‌گردد. فرض کنید $\mathbf{W} \in \mathbb{R}^{d_{\text{in}} \times d_{\text{out}}}$ ماتریس وزن هدف باشد که توسط تقریب کم‌مرتبه یا پراکنده با رتبه موثر $r$ جایگزین شده است. کران بالای خطای بازسازی بر مبنای نرم فروبنیوس به صورت زیر تصریح می‌شود:
\begin{equation}
\label{eq:spectral_bound_wgqa_attention}
\| \mathbf{W} - \mathbf{W}_{r} \|_F = \sqrt{\sum_{i=r+1}^{\min(d_{\text{in}}, d_{\text{out}})} \sigma_i^2(\mathbf{W})} \le \epsilon \| \mathbf{W} \|_F
\end{equation}
که در آن $\sigma_i(\mathbf{W})$ بیانگر $i$-امین مقدار منفرد مرتب‌شده ماتریس به صورت نزولی است. نسبت حفظ انرژی طیفی $\eta_r = \frac{\sum_{i=1}^r \sigma_i^2}{\sum_{i=1}^{d} \sigma_i^2}$ تضمین می‌کند که با انتخاب آستانه $\tau \ge 0.90$، بیش از نود درصد از پراکندگی واریانس داده‌ها حفظ شده و افت سیگنال در انتشار گرادیان رو به عقب به حداقل ممکن کاهش می‌یابد.

\paragraph{تحلیل مجانبی و پیچیدگی اشتراک حافظه سخت‌افزاری}
از منظر پیچیدگی محاسباتی و انتقال داده‌ها در سلسومراتب حافظه (حافظه سراسری \lr{HBM} به حافظه محلی \lr{SRAM})، عملیات ضرب ماتریسی در این مکانیزم نیازمند تخصیص بافر موقت با ابعاد فاکتورگیری‌شده است. برای پردازش دنباله‌ای به طول $N$ با بعد مخفی $D$، بار محاسباتی هسته از فرم درجه دو به رابطه زیر تقلیل می‌یابد:
\begin{equation}
\label{eq:complexity_wgqa_attention}
\mathcal{C}_{\text{wgqa\_attention}} = \mathcal{O}\left( B \cdot N \cdot D \cdot r + B \cdot H \cdot N \cdot d_k \right) \ll \mathcal{O}\left( B \cdot H \cdot N^2 \cdot d_k \right)
\end{equation}
این بهبود مرتبه پیچیدگی، با کاهش چشمگیر حجم تبادل تانسورها در گذرگاه‌های ارتباطی کارت گرافیک، مانع از اشباع پهنای باند حافظه (پدیده \lr{Memory-Bound}) در فرآیند استنتاج طولانی‌مدت می‌گردد.
